
% ============================
% CHAPTERS 11--15 FOR THE TOME
% Cosmology: White-State Cores and Cycles
% ============================

\part{Cosmology: White-State Cores and Cycles}

\chapter{Phase Cosmology and Background Solutions}

\section{Homogeneous and Isotropic Ansatz}

We consider a spatially homogeneous and isotropic cosmology with metric
\[
  ds^2 = -dt^2 + a(t)^2 \left( \frac{dr^2}{1-kr^2} + r^2 d\Omega^2 \right),
\]
where $a(t)$ is the scale factor and $k=0,\pm 1$ the spatial curvature index.

The phase field and chromatic fields are taken to be spatially homogeneous:
\[
  \tau = \tau(t), \qquad \chi_a = \chi_a(t),
\]
while the ER network contributes an effective homogeneous energy component $\rho_{ER}(t)$.

\section{Effective Friedmann Equations}

From the MK--43 Ultra action we obtain modified Friedmann equations of the form
\begin{align}
  H^2 &=
  \frac{8\pi G}{3}\left(\rho_\tau + \rho_\chi + \rho_{ER} + \rho_{\mathrm{m}}\right)
  - \frac{k}{a^2},
  \\
  \frac{\ddot{a}}{a} &=
  -\frac{4\pi G}{3}\left(
    \rho_\tau + \rho_\chi + \rho_{ER} + \rho_{\mathrm{m}}
    + 3p_\tau + 3p_\chi + 3p_{ER} + 3p_{\mathrm{m}}
  \right),
\end{align}
where $\rho_\tau$, $\rho_\chi$, and $\rho_{ER}$ denote energy densities in the phase, chromatic, and ER sectors, respectively, and $p_\tau$, $p_\chi$, $p_{ER}$ are the associated pressures.

\section{Role of the Phase Field in Expansion}

The phase field contributes an effective equation of state $w_\tau = p_\tau/\rho_\tau$ which can interpolate between stiff, radiation-like, or dark-energy-like behavior depending on the potential $V(\tau)$ and the magnitude of $\dot{\tau}$.
In particular, quasi-de Sitter phases can arise from regions where $V(\tau)$ dominates and $\dot{\tau}$ is small.

\chapter{Cosmology with Chromatic Sectors}

\section{Chromatic Energy Densities}

The chromatic fields contribute separate energy densities
\[
  \rho_\chi =
  \frac{1}{2} G^{ab} \dot{\chi}_a \dot{\chi}_b + U(\bm{\chi}),
\]
with pressures determined by the kinetic and potential balance.
Different chromatic modes can redshift at different rates, leading to distinct cosmological eras dominated by different bands.

\section{Chromatic Era Transitions}

Transitions between chromatic-dominated eras correspond to changes in which $\chi_a$ carries the dominant fraction of $\rho_\chi$.
Such transitions can imprint features in the expansion history and may leave signatures in primordial perturbations.

\section{Approach to the White-State}

As the universe contracts or energy becomes highly compressed, the chromatic amplitude vector $\bm{\chi}$ tends toward the symmetric configuration defining the white-state core.
This behaviour is central to nonsingular cosmological cycles in MK--43 Ultra.

\chapter{White-State Cores in Cosmological Bounces}

\section{Bounce Mechanism}

In a contracting phase, the total energy density
\[
  \rho_{\mathrm{tot}} = \rho_\tau + \rho_\chi + \rho_{ER} + \rho_{\mathrm{m}}
\]
approaches a finite critical value $\rho_{\mathrm{crit}}$ as the white-state is approached.
Rather than diverging, geometric and chromatic terms conspire to produce a minimal scale factor $a_{\mathrm{min}}$ at which
\[
  \dot{a} = 0, \qquad \ddot{a} > 0,
\]
implementing a bounce.

\section{White-State as Nonsingular Attractor}

The white-state core $\Wcore$ acts as a nonsingular attractor configuration in the space of cosmological solutions.
Nearby trajectories in the phase--chromatic--ER configuration space are drawn into a neighbourhood of $\Wcore$ before re-expanding along new directions in field space after the bounce.

\section{Entropy and Information Across the Bounce}

Information associated with pre-bounce structures is not necessarily erased.
Instead, it can be encoded in:

\begin{itemize}
  \item the ER network pattern at the bounce,
  \item residual phase correlations,
  \item the distribution of near-white chromatic perturbations.
\end{itemize}

Thus MK--43 Ultra allows for partial memory between cycles without requiring a singularity.

\chapter{ER Foam Evolution Through Cosmological Cycles}

\section{Coarse-Grained ER Density}

We introduce a coarse-grained ER density parameter $\rho_{ER}^{\mathrm{(eff)}}$ encoding the average strength and connectivity of the ER network at cosmological scales.
Its evolution equation may schematically be written as
\[
  \dot{\rho}_{ER}^{\mathrm{(eff)}}
  =
  \mathcal{F}\big(\rho_{ER}^{\mathrm{(eff)}}, \rho_\chi, \rho_\tau, a(t)\big),
\]
for some model-dependent functional $\mathcal{F}$.

\section{Network Reconfiguration at the Bounce}

Near the bounce, microscopic ER links are strongly reconfigured as chromatic fields approach the white-state.
This leads to a ``rewiring'' of the ER foam which seeds the adjacency structure of the next cycle.
The post-bounce ER network is thus statistically related but not identical to the pre-bounce network.

\section{Large-Scale Locality Emergence}

As the universe re-expands and the ER connectivity dilutes, an approximately local spacetime structure re-emerges on large scales.
Effective locality is therefore a derived property of the ER foam state in an expanding phase.

\chapter{Perturbations and Observational Windows}

\section{Linear Perturbations Around Cosmological Backgrounds}

We consider small perturbations around a homogeneous background:
\[
  \tau(t,\vec{x}) = \tau_0(t) + \delta\tau(t,\vec{x}), \qquad
  \chi_a(t,\vec{x}) = \chi_{a,0}(t) + \delta\chi_a(t,\vec{x}),
\]
with metric perturbations in a suitable gauge.
The coupled evolution equations for $(\delta\tau,\delta\chi_a,\delta g_{\mu\nu})$ determine the spectrum of cosmological perturbations.

\section{Signatures of Chromatic and Phase Dynamics}

Deviations from standard $\Lambda$CDM predictions may manifest as:

\begin{itemize}
  \item modified primordial power spectra,
  \item scale-dependent features from chromatic transitions,
  \item non-Gaussianities linked to ER network fluctuations,
  \item possible imprints of the bounce in low-$\ell$ observables.
\end{itemize}

\section{Future Directions}

A systematic confrontation of MK--43 Ultra cosmology with data would require:

\begin{itemize}
  \item specifying a concrete potential $V(\tau)$ and chromatic potential $U(\bm{\chi})$,
  \item modeling $\rho_{ER}^{\mathrm{(eff)}}$ explicitly,
  \item evolving perturbations numerically through the bounce,
  \item comparing predictions to CMB, large-scale structure, and gravitational wave observations.
\end{itemize}
These tasks are left for future work; this chapter provides the conceptual and formal scaffolding needed for such studies.
